\SourceLeaf{276}%
\noindent Wendet man diese Reihe auf die oben betrachteten Einzelfälle an, so erhält man
\[
\begin{aligned}
\text{für}\quad p&=-6,\quad\log m=+0{,}000001050448\\
\text{für}\quad p&=+6,\quad\log m=-0{,}000001096531
\end{aligned}
\]
Die endliche Formel 4, die man auch so schreiben kann,
\[
\begin{aligned}
m&=\frac{\alpha A\cos(Q+q)\BareRoot{(1-ee\sin(P+p)^2)}}{a\cos(P+p)}\\
&=\frac{\cos\eta\cos(Q+q)\BareRoot{(1-ee\sin(P+p)^2)}}{\cos\theta\cos(P+p)}
\end{aligned}
\]
ergibt bei einer Rechnung mit zehnstelligen Logarithmen bis auf die zehnte Ziffer genau dasselbe.
\Article{8}
\indent Für die umgekehrte Aufgabe, bei der \(q\) gegeben und \(p\) gesucht ist, ist die Entwicklung in eine Reihe noch wichtiger, da die endliche Formel 3 hier nur auf indirektem Weg zum Ziel führen könnte. Der Satz von \textsc{Taylor} ergibt
\[
p=\frac{\mathrm{d}w}{\mathrm{d}U}\mathord{.}q-\frac{\mathrm{dd}w}{2\mathrm{d}U^2}\mathord{.}qq+\frac{\mathrm{d}^3w}{6\mathrm{d}U^3}\mathord{.}q^3-\text{und so weiter}
\]
wobei für die Ableitungen diejenigen bestimmten Werte einzusetzen sind, die zu \(q=0\) oder zu \(U=90^0-Q\), \(w=90^0-P\) gehören. Für die allgemeinen Werte der ersten drei Ableitungen ergeben sich die folgenden Ausdrücke.
\[
\frac{\mathrm{d}w}{\mathrm{d}U}=\frac{(1-ee\cos w^2)\sin w}{\alpha(1-ee)\sin U}
\]
\[
\begin{aligned}
\frac{\mathrm{dd}w}{\mathrm{d}U^2}={}&\frac{(1-ee\cos w^2)\sin w}{\alpha\alpha(1-ee)^2\sin U^2}\\
&\quad(\alpha(1-ee)\cos U-\cos w+ee\cos w(\cos w^2-2\sin w^2))
\end{aligned}
\]
\[
\begin{aligned}
\frac{\mathrm{d}^3w}{\mathrm{d}U^3}={}&\frac{(1-ee\cos w^2)\sin w}{\alpha^3(1-ee)^3\sin U^3}\{\alpha\alpha(1-ee)^2(\cos U^2+2\sin U^2)\\
&-3\alpha(1-ee)\cos U\cos w(1-ee(\cos w^2-2\sin w^2))\\
&+\cos w^2-\sin w^2-ee(2\cos w^4-12\cos w^2\sin w^2+2\sin w^4)\\
&+e^4(\cos w^6-11\cos w^4\sin w^2+6\cos w^2\sin w^4)\}
\end{aligned}
\]
\indent Die beiden folgenden Koeffizienten habe ich ebenfalls vollständig entwickelt. Um Platz zu sparen, gebe ich sie hier nicht an, da sie nur Zwischengrößen auf dem Weg zu den Endergebnissen sind. Die Endergebnisse erhält man nach der %
